% !TeX root = ../../Book4.tex
% 纲要标题：## 第10章　扩张、乘积与万有系数
% 纲要标签：b4:ch:09
% 纲要位置：计划纲要.md，第 3435 行。

\chapter{扩张、乘积与万有系数}
\label{b4:ch:09}
\BookChapterNumber{b4:ch:09}

\begin{chapterabstract}
本章把 Ext 与模扩张联系起来，研究扩张的乘积，并证明主理想整环上的万有系数定理与 Künneth 公式，区分自然正合列和非典范分裂。
\end{chapterabstract}

\begin{learninggoals}
\begin{itemize}
\item 从投射消解构造扩张类，并用 Baer 和与 Yoneda 乘积计算它们。
\item 说明高阶扩张等价为何允许箭头方向交替，区分扩张类与中间对象的同构类型。
\item 核对张量微分的符号，辨别平坦性及有界性在保持拟同构中的作用。
\item 应用万有系数与 Künneth 公式，保留附加的 Ext、Tor 项并辨别分裂的自然性。
\end{itemize}
\end{learninggoals}

考察以 \(\mathbb Z/2\mathbb Z\) 为两端的短正合列。若中间项也是两个二元群的直积，可以从右端选一个同态截面；若中间项是 \(\mathbb Z/4\mathbb Z\)，右端的非零元却没有阶为 \(2\) 的原像，截面便不存在。两条序列的端点相同，中间的接合方式却不同。Ext 的一阶群把这些接合方式组织起来；把两条短正合列首尾相接，又提示高阶 Ext 中的乘积。另一方面，对复形取张量后，同调可能新增 Tor 项。我们将由这些可见的差别走向扩张积、万有系数与 Künneth 公式。

\ProseParagraphBreak
扩张的拼接与复形的张量积处理的是不同数据。前者把扩张次数相加；后者先构造一个新复形，再考察其同调。\cref{b4:tab:extension-tensor-comparison}将这两种操作与它们涉及的同调比较映射并列。

\begin{table}[htbp]
\centering
\caption{扩张乘积、张量复形与同调比较}
\label{b4:tab:extension-tensor-comparison}
\begin{tabularx}{\textwidth}{@{}p{0.19\textwidth}>{\raggedright\arraybackslash}X>{\raggedright\arraybackslash}X@{}}
\toprule
构造 & 输入的数据 & 所得对象或映射 \\
\midrule
Yoneda 乘积
& \(\alpha\in\operatorname{Ext}_R^m(N,L)\)，\newline \(\beta\in\operatorname{Ext}_R^n(M,N)\)
& 拼接得到 \(\alpha\beta\in\operatorname{Ext}_R^{m+n}(M,L)\) \\
张量总复形
& 右 \(R\)-模链复形 \(C\) 与左 \(R\)-模链复形 \(D\)
& 复形 \(C\otimes_R D\)，微分包含次数符号 \((-1)^p\) \\
同调比较映射
& 两个复形的同调类的张量积，\(p+q=n\)
& \(\kappa_n\) 映到 \(H_n(C\otimes_R D)\)；Künneth 公式计算其余核 \\
\bottomrule
\end{tabularx}
\end{table}


本章使用\chapref{b4:ch:05}、\chapref{b4:ch:06}的复形和总化，以及\chapref{b4:ch:07}、\chapref{b4:ch:08}的消解、Ext、Tor 和维数平移。\secref{b4:sec:0901}以左模为对象；张量部分逐项说明模的侧别，万有系数和 Künneth 主定理则在主理想整环上讨论。为方便同调公式，后三节采用链下标 \(C_n=C^{-n}\)，不改变本卷上链微分与移位的约定。第三卷的 Koszul 复形提供了另一种计算方式，可与这里的公式对照。

\ProseParagraphBreak
扩张与乘积可参阅 Weibel 的《An Introduction to Homological Algebra》\cite[Section 3.4]{Weibel1994HomologicalAlgebra}；系数变换与复形张量积的公式见\cite[Section 3.6]{Weibel1994HomologicalAlgebra}。本章对高阶扩张给出具体的双向构造，对主理想整环上的分裂逐步证明其存在，并用链自同构说明为何不能把它称为自然分裂。

% !TeX root = ../../Book4.tex
% 纲要标题：### 10.1 Yoneda 扩张
% 纲要标签：b4:sec:0901
% 纲要位置：计划纲要.md，第 3437 行。

\section{Yoneda 扩张}
\label{b4:sec:0901}

扩张不仅包含中间模，还包含核与商的指定映射。把扩张粘接起来得到的乘积，需要与消解中提升复形映射的乘积一致。

% 纲要内容：
%
% * Ext^1 与短正合扩张
% * 高阶扩张及等价
% * Yoneda 乘积与消解乘积的一致性
%

\subsection{\texorpdfstring{\(\operatorname{Ext}^1\) 与短正合扩张}{Ext1 与短正合扩张}}
\label{b4:subsec:0901:01}

给定左 \(R\)-模 \(M,N\)，\(\operatorname{Ext}_R^1(M,N)\) 既可由消解计算，也能描述把 \(N\) 放在核位置、把 \(M\) 放在商位置的所有短正合列。这里分类的数据包括两端的指定映射；仅知道中间模的同构类型还不够。

\begin{definition}\label{b4:def:yoneda-extension}
设 \(R\) 为含幺环，\(M,N\) 为幺性左 \(R\) 模。短正合列
\[
 \xi:\quad 0\longrightarrow N\xrightarrow{i}E\xrightarrow{p}M\longrightarrow0
\]
称为 \(M\) 被 \(N\) 的一个\term[kuozhang]{扩张}[extension]。给定另一个扩张 \(\xi':0\to N\xrightarrow{i'}E'\xrightarrow{p'}M\to0\)。若存在 \(R\) 模同态 \(f:E\to E'\)，满足 \(fi=i'\)、\(p'f=p\)，则称两扩张等价。这里两端的映射固定为恒等映射；两行短正合且两端为同构，由\cref{b4:thm:five-lemma}应用于包含两端零对象的五项序列，可知 \(f\) 自动为同构。仅有 \(E\cong E'\) 还不够，所用同构必须同时保持指定的核嵌入与商映射。
\end{definition}

\begin{theorem}\label{b4:thm:ext1-extensions}
设 \(R\) 为含幺环，\(M,N\) 为幺性左 \(R\) 模。\(M\) 被 \(N\) 的短正合扩张的等价类与 \(\operatorname{Ext}_R^1(M,N)\) 自然一一对应，分裂扩张对应零元。
\end{theorem}
\begin{proof}
选定正合列 \(0\to K\xrightarrow{j}P\xrightarrow{\varepsilon}M\to0\)，其中 \(P\) 投射。\cref{b4:thm:ext-long-exact}在第一变量给出的低次片段为
\[
 \operatorname{Hom}_R(P,N)\xrightarrow{j^*}\operatorname{Hom}_R(K,N)
 \longrightarrow\operatorname{Ext}_R^1(M,N)
 \longrightarrow\operatorname{Ext}_R^1(P,N)=0.
\]
最后一项由 \(P\) 的投射性消失，且 \(j^*(a)=aj\)，所以
\[
 \operatorname{Ext}_R^1(M,N)
 \cong\operatorname{Hom}_R(K,N)\big/\{\,a j:a\in\operatorname{Hom}_R(P,N)\,\}.
\]
分母恰好由可以从 \(K\) 延拓到 \(P\) 的映射组成。对扩张 \(\xi\)，由 \(P\) 投射且 \(p:E\to M\) 满射，可将 \(\varepsilon\) 提升为 \(t:P\to E\)，使 \(pt=\varepsilon\)。随后 \(ptj=\varepsilon j=0\)，故核映射 \(i:N\to E\) 的泛性质给出唯一的 \(u:K\to N\)，使 \(iu=tj\)。提升得到 \(t\)，经过核才得到 \(u\)；后者记录了 \(P\) 中的关系在扩张里落入 \(N\) 的方式。换成另一提升 \(t'\)，有 \(p(t'-t)=0\)，因而 \(t'-t=ia\)（\(a:P\to N\)），再由 \(i\) 单射得 \(u'-u=aj\)。扩张态射与 \(i,p\) 交换，将一套提升送到另一扩张的提升，故也保持这个余类。

反过来，希望把 \(P\) 中的 \(j(k)\) 与 \(N\) 中的 \(u(k)\) 识别，即令 \([(0,j(k))]=[(u(k),0)]\)。把两边相减便得到应商掉的关系 \((u(k),-j(k))\)，所以沿 \(u:K\to N\) 作推出得到
\[
 E_u=(N\oplus P)/\{\,(u(k),-j(k)):k\in K\,\}.
\]
由 \(n\mapsto[(n,0)]\) 及 \([(n,p)]\mapsto\varepsilon(p)\) 得到扩张。第一个映射单射，因为 \(j\) 单射；第二个映射的核由 \(\varepsilon(p)=0\) 即 \(p=j(k)\) 算出，此时 \([(n,j(k))]=[(n+u(k),0)]\)，故核恰好是 \(N\) 的像。如果 \(u-v=a j\)，考虑映射 \([(n,p)]\mapsto[(n+a(p),p)]\)。它把被商掉的关系送到
\[
 (u(k),-j(k))\longmapsto(u(k)-aj(k),-j(k))=(v(k),-j(k)),
\]
因而良定义；用 \(-a\) 得到逆映射，且两端映射保持不变。对原来的 \(\xi\)，映射 \([(n,p)]\mapsto i(n)+t(p)\) 也良定义，因为 \(iu=tj\)，并给出 \(E_u\to E\) 的扩张态射，因而是同构。两个构造互逆。

若类为零，则 \(u=aj\)（\(a:P\to N\)）。令 \(t'=t-ia\)，就有 \(pt'=\varepsilon\) 及
\[
 t'j=tj-iaj=iu-iu=0.
\]
因此 \(t'\) 唯一经过余核 \(\varepsilon:P\to M\)，写为 \(s\varepsilon\)；由 \(ps\varepsilon=\varepsilon\) 且 \(\varepsilon\) 满射，得到 \(ps=1_M\)。反之，若有截面 \(s\)，可取提升 \(t=s\varepsilon\)，它在 \(K\) 上为零，故对应零类。对 \(N\to N'\) 作推出相当于把 \(u\) 后复合，对 \(M'\to M\) 作拉回相当于提升该映射到投射消解后预复合，故对应自然。
\end{proof}

扩张类的加法也有直接构造。对两个扩张先取直和，再沿对角映射 \(M\to M\oplus M\) 拉回，使两个中间元素具有同一个商像；所得核仍有两份 \(N\)。随后沿加法 \(N\oplus N\to N\) 推出，把核元素 \((n,n')\) 合并为 \(n+n'\)。所得扩张称为\term[baerhe]{Baer 和}[Baer sum]。\footnote{拜尔（Reinhold Baer，1902-1979），德裔数学家，研究群的扩张、可解群和幂零群。}加法的核由 \((n,-n)\) 组成，故中间模具体为
\[
 (E\times_M E')/\{\,(i(n),-i'(n)):n\in N\,\}.
\]

\begin{proposition}\label{b4:prop:baer-addition}
设 \(R\) 为含幺环，\(M,N\) 为幺性左 \(R\) 模。Baer 和在扩张等价类上给出交换群结构；\cref{b4:thm:ext1-extensions}的对应是群同构。
\end{proposition}
\begin{proof}
从同一个 \(P\to M\) 分别提升到 \(E,E'\)，因两套提升的商像相同，可合并为到纤维积的提升；再作推出后，在 \(K\) 上诱导的映射就是 \(u+u'\)。因而 Baer 和对应 \(\operatorname{Ext}^1\) 中的加法。前一定理给出的双射于是把商群中已成立的结合律、交换律及零元、负元传到扩张类上，也证明结果不依赖代表元。具体地，分裂扩张为零，沿 \(-1_N\) 推出得到负元。
\end{proof}

例如，对素数 \(p\)，\(\operatorname{Ext}_{\mathbb Z}^1(\mathbb Z/p\mathbb Z,\mathbb Z/p\mathbb Z)=\mathbb Z/p\mathbb Z\)。非零类的中间群都同构于 \(\mathbb Z/p^2\mathbb Z\)，但指定的核嵌入可以不同，所以这些扩张仍有 \(p-1\) 个不同的非零类。这个区别在用扩张恢复结构时不能省略。

\subsection{高阶扩张及等价}
\label{b4:subsec:0901:02}

\begin{definition}\label{b4:def:higher-yoneda-extension}
设 \(R\) 为含幺环，\(M,N\) 为幺性左 \(R\) 模，\(n\ge1\) 为整数。给定左 \(R\) 模正合列
\[
 \xi:\quad 0\longrightarrow N\longrightarrow E_{n-1}\longrightarrow\cdots
 \longrightarrow E_0\longrightarrow M\longrightarrow0
\]
称为一个 \(n\) 阶扩张。两个这样的扩张之间，两端为恒等映射的链映射称为扩张态射。由扩张态射生成的等价关系称为 \term[yonedakuozhang]{Yoneda 扩张}[Yoneda extension]的等价关系：允许有限串扩张态射，箭头方向可以交替。这些扩张的等价类族暂记为 \(\operatorname{YExt}_R^n(M,N)\)；下面的定理将证明它可由一个集合参数化。
\end{definition}
\symidx{\(\operatorname{YExt}_R^n(M,N)\)：\(n\) 阶 Yoneda 扩张的等价类}

当 \(n>1\) 时，中间各项的映射不必可逆。例如令 \(A=\mathbb Z/2\mathbb Z\)，考察两个整数模二阶扩张
\[
\begin{aligned}
0&\longrightarrow A\xrightarrow{a\mapsto(a,0)}A\oplus\mathbb Z
 \xrightarrow{(a,z)\mapsto2z}\mathbb Z
 \xrightarrow{z\mapsto\bar z}A\longrightarrow0,\\
0&\longrightarrow A\xrightarrow{1_A}A\xrightarrow{0}A
 \xrightarrow{1_A}A\longrightarrow0.
\end{aligned}
\]
第一列中间的核为 \(A\oplus0\)，像为 \(2\mathbb Z\)，所以两列均正合。中间分别取投影 \((a,z)\mapsto a\) 和取模映射 \(z\mapsto\bar z\)，两端取恒等，就得到从第一列到第二列的扩张态射：中间方块两条复合均为零，其余方块也交换。中间项却分别是无限群与有限群，不可能同构。高阶等价因而不能要求逐项同构；要由这些不必可逆的态射得到等价关系，还须允许有限箭头串及箭头反向。下面同时说明这些类确实构成集合，并给出计算它们的方法。

\begin{theorem}\label{b4:thm:yoneda-ext-comparison}
设 \(R\) 为含幺环，\(M,N\) 为幺性左 \(R\) 模，\(n\ge1\) 为整数。扩张等价类可组成集合，并有对 \(M\) 反变、对 \(N\) 协变的自然双射
\[
 \operatorname{YExt}_R^n(M,N)\simeq\operatorname{Ext}_R^n(M,N).
\]
沿 \(M\) 的对角映射拉回、沿 \(N\) 的加法映射推出所定义的高阶 Baer 和，对应右侧的加法。
\end{theorem}
\begin{proof}
固定投射消解 \(P_\bullet\xrightarrow{\varepsilon}M\)，令 \(\Omega^nM=\ker(P_{n-1}\to P_{n-2})\)，其中 \(P_{-1}=M\)，\(P_0\to P_{-1}\) 为增广。记由微分诱导的满射为 \(\pi_n:P_n\to\Omega^nM\)。\cref{b4:thm:dimension-shift-basic,b4:prop:iterated-dimension-shift}的平移公式及短正合列 \(0\to\Omega^nM\to P_{n-1}\to\Omega^{n-1}M\to0\) 给出
\[
 \operatorname{Ext}_R^n(M,N)
 \simeq\operatorname{Hom}_R(\Omega^nM,N)
 /\operatorname{res}\operatorname{Hom}_R(P_{n-1},N).
\]
也可直接从消解读出此式。例如 \(n=2\) 时，\(\Omega^2M=\ker\mathrm{d}_{P,1}=\operatorname{im}\mathrm{d}_{P,2}\)。次数二的余圈 \(b:P_2\to N\) 满足 \(b\mathrm{d}_{P,3}=0\)，而 \(\ker\pi_2=\operatorname{im}\mathrm{d}_{P,3}\)，所以唯一写成 \(b=u\pi_2\)。次数二的余边是 \(a\mathrm{d}_{P,2}\)（\(a:P_1\to N\)），其下降映射恰为 \(a|_{\Omega^2M}\)。任意 \(n\ge1\) 同样有 \(\ker\pi_n=\operatorname{im}\mathrm{d}_{P,n+1}\)，因此余圈唯一经过 \(\Omega^nM\)，余边则对应商公式分母中的限制映射。这里沿用\secref{b4:sec:0802}投射消解计算中的无符号微分 \(f\mapsto f\mathrm{d}_P\)。

给定 \(n\) 阶扩张，将 \(E_0\to M\) 记为 \(\mathrm{d}_{E,0}\)，并将 \(P_0\to M\) 记为 \(\mathrm{d}_{P,0}\)。依次利用 \(P_i\) 的投射性，构造与微分、增广交换的提升 \(t_i:P_i\to E_i\)，\(0\le i<n\)。最后 \(t_{n-1}|_{\Omega^nM}\) 被下一微分零化，故经最左端的核嵌入 \(i:N\to E_{n-1}\) 分解，得到 \(u:\Omega^nM\to N\)。

为证明其余类与提升无关，只需把两套提升在最左端的差写成一个 \(P_{n-1}\to N\) 映射的限制。不能直接断言这一差落入 \(N\)：较低次数的提升也可能不同，须先逐层消去由它们传来的误差。对两套提升 \(t,s\)，从最低次数起构造 \(h_i:P_i\to E_{i+1}\)，使
\[
 t_i-s_i=\mathrm{d}_Eh_i+h_{i-1}\mathrm{d}_P\qquad(0\le i<n-1),\qquad h_{-1}=0.
\]
零次的差被增广零化。若已构造前一步，令
\[
 e_i=t_i-s_i-h_{i-1}\mathrm{d}_{P,i}.
\]
由两套提升的交换等式及上一同伦等式，得
\[
\begin{aligned}
\mathrm{d}_{E,i}e_i
&=(t_{i-1}-s_{i-1})\mathrm{d}_{P,i}
  -\mathrm{d}_{E,i}h_{i-1}\mathrm{d}_{P,i}\\
&=h_{i-2}\mathrm{d}_{P,i-1}\mathrm{d}_{P,i}=0
\qquad(i\ge1).
\end{aligned}
\]
所以尚未消去的误差落在核中；中间正合性把这个核识别为 \(E_{i+1}\) 的像，\(P_i\) 投射才使它提升为 \(h_i\)。到 \(i=n-1\) 时，核已经是 \(i(N)\)，于是存在 \(a:P_{n-1}\to N\) 使
\[
 t_{n-1}-s_{n-1}-h_{n-2}\mathrm{d}_{P,n-1}=ia.
\]
限制到 \(\Omega^nM\) 后，微分项为零，再由 \(i\) 单射得到 \(u_t-u_s=a|_{\Omega^nM}\)，这恰在商公式的分母中。\(n=1\) 时直接用 \(t_0-s_0\) 即可。扩张态射把一套提升变成另一扩张的提升，因而整个箭头串都保持该余类。

对任意 \(u:\Omega^nM\to N\)，把截短消解
\[
 0\longrightarrow\Omega^nM\longrightarrow P_{n-1}\longrightarrow\cdots
 \longrightarrow P_0\longrightarrow M\longrightarrow0
\]
的最左端沿 \(u\) 推出。新项是
\(E_{n-1}(u)=(N\oplus P_{n-1})/\{(u(x),-x)\}\)，其余项保持 \(P_i\)。与一阶情形相同，核及像的计算证明新列正合。若 \(u-v=a|_{\Omega^nM}\)，在新项上用 \([(z,p)]\mapsto[(z+a(p),p)]\)，在其他项上用恒等映射，就得到两个新扩张的同构。

从任意原扩张取出的 \(u\) 又给出这个推出扩张到原扩张的态射：最左新项用 \([(z,p)]\mapsto i(z)+t_{n-1}(p)\)，其他项用 \(t_i\)。标准推出扩张的这些项取自 \(P_\bullet\)，原扩张的 \(E_i\) 却不必投射，所构造的中间箭头也不必可逆；它们之所以足以识别两扩张，正是因为高阶等价由扩张态射生成。因此两个方向互逆。特别地，每个扩张类都有一个由 \(u\in\operatorname{Hom}_R(\Omega^nM,N)\) 得到的代表；这个 Hom 集的上述商集便给出全部扩张类的集合参数化。

接着核对两变量的作用。若 \(v:N\to N'\)，沿 \(v\) 推出扩张，并把原来的提升复合到推出列中，最左端诱导的映射变成 \(vu:\Omega^nM\to N'\)。它对应消解余圈的后复合，故与 \(\operatorname{Ext}_R^n(M,v)\) 相容。

若 \(f:M'\to M\)，取投射消解 \(P'_\bullet\to M'\)，由\cref{b4:thm:resolution-comparison}选取提升 \(f\) 的比较映射 \(c:P'_\bullet\to P_\bullet\)。记
\[
 c_\Omega=c_{n-1}|_{\Omega^nM'}:\Omega^nM'\longrightarrow\Omega^nM.
\]
这里 \(f\) 决定要拉回的扩张，\(c\) 实现 Ext 计算中的预复合，\(c_\Omega\) 则记录它在关系模上的作用。拉回列的最低项为 \(E_0\times_M M'\)，其提升明确取为
\[
 P'_0\longrightarrow E_0\times_M M',\qquad
 x\longmapsto(t_0c_0(x),\varepsilon_{P'}(x));
\]
它落在纤维积中，因为 \(\mathrm{d}_{E,0}t_0c_0=f\varepsilon_{P'}\)。其余次数用原提升与 \(c_i\) 的复合，最左端映射因而成为 \(u c_\Omega\)。若余圈 \(b:P_n\to N\) 写为 \(b=u\pi_n\)，记相应满射为 \(\pi'_n:P'_n\to\Omega^nM'\)，则
\[
 b c_n=u\pi_nc_n=u c_\Omega\pi'_n,
\]
所以这个商类恰是 \(\operatorname{Ext}_R^n(f,N)\) 所给的类。若换另一比较提升 \(c'\)，由\cref{b4:prop:comparison-homotopy-unique}，存在同伦 \(h_i:P'_i\to P_{i+1}\)，使
\[
 c_i-c'_i=\mathrm{d}_{P,i+1}h_i+h_{i-1}\mathrm{d}_{P',i},\qquad h_{-1}=0.
\]
在 \(\Omega^nM'\) 上第二项为零，因此
\[
 u c_\Omega-u c'_\Omega
 =\left.(u\pi_nh_{n-1})\right|_{\Omega^nM'}.
\]
其中 \(h_{n-1}:P'_{n-1}\to P_n\)，故 \(u\pi_nh_{n-1}\) 定义在整个 \(P'_{n-1}\) 上，右侧恰属于商公式的分母。这既证明拉回与 Ext 的预复合相容，也证明更换比较提升不改变所得类。

最后把换消解与换扩张提升区分开来。若 \(P'_\bullet\) 也是 \(M\) 的投射消解，上面的构造取 \(f=1_M\)，就把两种商描述按 \([u]\mapsto[u c_\Omega]\) 比较。取反向比较后，两种复合分别与相应的恒等映射同伦，所以诱导互逆的同构；余圈等式也表明这正是\chapref{b4:ch:08}用于识别两种 Ext 计算的比较同构。因而不同消解给出的识别在该同构下相容。前面固定消解时已经证明换扩张中的提升、沿扩张态射改变代表均不改变商类，这里消去的则是投射消解本身的选择。

对扩张 \(\xi,\xi'\)，先逐项取直和 \(\xi\oplus\xi'\)，得到右端为 \(M\oplus M\)、左端为 \(N\oplus N\) 的列；再在右端沿 \(\Delta_M:M\to M\oplus M\) 拉回，最后在左端沿 \(+_N:N\oplus N\to N\) 推出。也就是
\[
 \xi+\xi'=(+_N)_*\Delta_M^*(\xi\oplus\xi').
\]
分别取 \(P_\bullet\) 到两列的提升，再用逐项对角映射 \(P_\bullet\to P_\bullet\oplus P_\bullet\) 提升到拉回列。在最左端，所得映射是 \((u,u'):\Omega^nM\to N\oplus N\)；推出后成为 \(+_N(u,u')=u+u'\)。所以 Baer 和对应商集中的加法。特别地，它不依赖代表扩张，并由此给出与 Ext 相同的交换群结构。
\end{proof}

对 \(n=0\)，约定 \(\operatorname{YExt}_R^0(M,N)=\operatorname{Hom}_R(M,N)\)。在没有足够投射或内射对象的一般交换范畴中，仍可讨论扩张及其箭头串；本节已证的与导出函子的比较则是在模范畴中进行的。

\subsection{Yoneda 乘积与消解乘积}
\label{b4:subsec:0901:03}

若一列以 \(N\) 为商，另一列以 \(N\) 为核，就能把它们首尾连接。例如两个一次扩张
\[
0\longrightarrow L\xrightarrow{i}E\xrightarrow{v}N\longrightarrow0,
\qquad
0\longrightarrow N\xrightarrow{j}F\xrightarrow{p}M\longrightarrow0
\]
拼接为
\[
0\longrightarrow L\xrightarrow{i}E\xrightarrow{jv}F
\xrightarrow{p}M\longrightarrow0.
\]
由于 \(j\) 单射，\(\ker(jv)=\ker v=\operatorname{im}i\)；由于 \(v\) 满射，\(\operatorname{im}(jv)=j(N)=\ker p\)，所以拼接列仍正合。中间箭头是明确的复合 \(E\to N\to F\)，而不是将两列中的 \(N\) 直接略去。这一操作给出 Ext 群之间并非来自普通数乘的乘法。

\begin{definition}\label{b4:def:yoneda-product}
设 \(R\) 为含幺环，\(L,N,M\) 为幺性左 \(R\) 模，\(m,n\ge1\) 为整数。对 \(\alpha\in\operatorname{YExt}_R^m(N,L)\)、\(\beta\in\operatorname{YExt}_R^n(M,N)\)，把代表扩张在共同的 \(N\) 处拼接，删除两处 \(N\)，以原来的满射到 \(N\) 与从 \(N\) 出发的单射之复合作为中间箭头。所得 \((m+n)\) 阶扩张类记为 \(\alpha\beta\)，称为\term[yonedachengji]{Yoneda 乘积}[Yoneda product]。约定 \(\operatorname{YExt}_R^0(X,Y)=\operatorname{Hom}_R(X,Y)\)：左因子为零次同态 \(N\to L\) 时沿该同态推出，右因子为零次同态 \(M\to N\) 时沿该同态拉回；两个因子均为零次时取同态复合。
\end{definition}
\symidx{\(\alpha\beta\)：Yoneda 扩张类的乘积}

\begin{proposition}\label{b4:prop:yoneda-resolution-product}
设 \(R\) 为含幺环，\(L,N,M\) 为幺性左 \(R\) 模。Yoneda 乘积良定义、双加性且结合，并在 Yoneda 扩张与 Ext 的自然对应下等于下述消解乘积。取投射消解 \(P_\bullet\to M\)、\(Q_\bullet\to N\)，以 \(\varepsilon_Q:Q_0\to N\) 表示增广，设 \(m,n\ge0\) 为整数。若 \(b:P_n\to N\)、\(a:Q_m\to L\) 分别是 \(\operatorname{Hom}_R(P_\bullet,N)\)、\(\operatorname{Hom}_R(Q_\bullet,L)\) 中的余圈，把 \(b\) 逐次提升为
\[
 t_i:P_{n+i}\longrightarrow Q_i,\qquad
 \varepsilon_Qt_0=b,\qquad \mathrm{d}_Qt_i=t_{i-1}\mathrm{d}_P\quad(i\ge1),
\]
则乘积由余圈 \(a t_m:P_{n+m}\to L\) 代表。
\end{proposition}
\begin{proof}
拼接在共同项处正合，因为像和核都由被删除的 \(N\) 描述。若一个因子沿某个扩张态射改变，把该态射与另一列的恒等映射拼接，仍得到扩张态射。因此拼接保持等价关系的每个生成关系，也就保持由正向或反向扩张态射组成的任意有限箭头串。零次因子的拉回或推出也把扩张态射送到扩张态射，故这些情形同样良定义。三个正次数因子依次拼接不改变项及其复合箭头。零次因子位于外端时，拉回或推出的泛性质使两种顺序相容；它位于中间时，拉回的投影与推出的结构映射给出两种拼接列之间的扩张态射。连续的零次因子则由同态复合的结合律处理。因而乘积结合。

对于消解公式，先沿满射 \(\varepsilon_Q:Q_0\to N\) 提升余圈 \(b:P_n\to N\)，得到 \(t_0:P_n\to Q_0\)。因此 \(Q\) 的零次已经对应 \(P\) 的第 \(n\) 次。由 \(b\mathrm{d}_{P,n+1}=0\)，\(t_0\mathrm{d}_{P,n+1}\) 落在 \(\ker\varepsilon_Q\)，可用 \(P_{n+1}\) 的投射性提升至 \(Q_1\)。以后每步同样先用 \(\mathrm{d}_P^2=0\) 核对落在核中，再提升；每向 \(Q\) 的高次推进一步，源也从 \(P_{n+i}\) 推进到 \(P_{n+i+1}\)。到第 \(m\) 步，恰可形成
\[
 P_{n+m}\xrightarrow{t_m}Q_m\xrightarrow{a}L,
 \qquad
 a t_m\mathrm{d}_{P,n+m+1}
 =a\mathrm{d}_{Q,m+1}t_{m+1}=0.
\]
这既说明所得是余圈，也说明乘积次数为何是 \(n+m\)。

对 \(m,n\ge1\)，要比较两种乘积，只须读取拼接扩张最左端的关系映射，证明其余圈为 \(at_m\)。用\cref{b4:thm:yoneda-ext-comparison}的推出构造代表 \(\alpha,\beta\)：把 \(\alpha\) 的最左新项记为 \(A\)，把 \(\beta\) 的最左新项记为 \(B\)，其余项来自两条消解。设 \(a,b\) 经合冲模的分解分别为
\[
 \bar a:\Omega^mN\to L,\qquad \bar b:\Omega^nM\to N,
\]
则这两个推出项为
\[
\begin{aligned}
 A&=(L\oplus Q_{m-1})/\{(\bar a(x),-x):x\in\Omega^mN\},\\
 B&=(N\oplus P_{n-1})/\{(\bar b(y),-y):y\in\Omega^nM\}.
\end{aligned}
\]
以 \(i_\alpha:L\to A\)、\(i_\beta:N\to B\) 表示单射，以 \(j_\alpha:Q_{m-1}\to A\)、\(j_\beta:P_{n-1}\to B\) 表示 \(q\mapsto[(0,q)]\) 型的映射。要将两段提升在共同的 \(N\) 处接起来，须使经 \(P_{n-1}\) 进入 \(B\) 与经 \(Q_0\to N\) 进入 \(B\) 的两条路径相同，所需等式正是下面的方块：
\[
\begin{tikzcd}[column sep=large,row sep=large]
 P_n \ar[r,"\mathrm{d}_{P,n}"] \ar[d,"t_0"']
   & P_{n-1} \ar[d,"j_\beta"] \\
 Q_0 \ar[r,"i_\beta\varepsilon_Q"'] & B.
\end{tikzcd}
\]
它交换，因为 \(j_\beta\mathrm{d}_{P,n}=i_\beta b=i_\beta\varepsilon_Qt_0\)。因此从 \(P_0\) 到 \(P_{n-2}\) 的恒等映射、\(P_{n-1}\to B\) 的 \(j_\beta\)，以及
\[
 P_{n+i}\xrightarrow{t_i}Q_i\quad(0\le i<m-1),\qquad
 P_{n+m-1}\xrightarrow{j_\alpha t_{m-1}}A
\]
组成到拼接列的逐项提升；当 \(m=1\) 或 \(n=1\) 时，相应的指标区间为空。若 \(m=1\)，拼接箭头 \(A\to B\) 由 \(i_\beta\varepsilon_Q:Q_0\to B\) 及 \(L\to B\) 的零映射在推出上诱导，所以上面的方块仍给出所需相容性。

这些映射给出到拼接列的提升以后，最左端的关系映射由 \(P_{n+m}\to P_{n+m-1}\to A\) 读出。下面的上方方块使用提升等式，下方方块使用 \(A\) 的推出关系，将这条路径与 \(P_{n+m}\to Q_m\to L\to A\) 比较：
\[
\begin{tikzcd}[column sep=large,row sep=large]
 P_{n+m} \ar[r,"\mathrm{d}_{P,n+m}"] \ar[d,"t_m"']
   & P_{n+m-1} \ar[d,"t_{m-1}"] \\
 Q_m \ar[r,"\mathrm{d}_{Q,m}"] \ar[d,"a"']
   & Q_{m-1} \ar[d,"j_\alpha"] \\
 L \ar[r,"i_\alpha"'] & A.
\end{tikzcd}
\]
下方方块交换，是因为推出关系给出 \(j_\alpha\mathrm{d}_{Q,m}=i_\alpha a\)。于是
\[
 j_\alpha t_{m-1}\mathrm{d}_{P,n+m}=i_\alpha a t_m.
\]
由于 \(\operatorname{im}\mathrm{d}_{P,n+m}=\Omega^{n+m}M\) 且 \(i_\alpha\) 单射，拼接列在最左端诱导的映射正是余圈 \(a t_m\) 经 \(\Omega^{n+m}M\) 的分解。因此拼接扩张对应所陈述的消解乘积。该对应也证明换 \(t_i\)、换余圈代表或换消解均不改变所得 Ext 类。

固定 \(b\) 的提升时，\((a+a')t_m=a t_m+a't_m\)；固定 \(a\) 时，两套 \(b,b'\) 的提升之和提升 \(b+b'\)。因此乘积双加性。若 \(m=0\)，写 \(a=h\varepsilon_Q\)，其中 \(h:N\to L\)，则 \(a t_0=h b\)，正是沿 \(h\) 推出的后复合。若 \(n=0\)，写 \(b=f\varepsilon_P\)，其中 \(f:M\to N\)，则 \(t_i\) 是提升 \(f\) 的比较映射，\(a t_m\) 正是自然性中沿 \(f\) 拉回的余圈。两个次数均为零时，这又化为同态复合。
\end{proof}

于是 \(\bigoplus_{n\ge0}\operatorname{Ext}_R^n(M,M)\) 成为带单位的分次环，单位是 \(1_M\)。这里没有一般的分次交换性断言：次数零部分已经可能是非交换环 \(\operatorname{End}_R(M)\)。例如在域 \(R=k\) 上取 \(M=k^2\)，这一部分就是矩阵环 \(M_2(k)\)。

\ProseParagraphBreak
即使次数零部分交换，正次数也未必满足分次交换律。取域 \(k\) 及交换环 \(R=k[\epsilon]/(\epsilon^2)\)，令 \(M=k=R/(\epsilon)\)。\secref{b4:sec:0701}已证明它有各项为 \(R\)、微分均为乘 \(\epsilon\) 的周期自由消解；\secref{b4:sec:0802}由施加 Hom 后微分为零得到每次 Ext 为 \(k\)。次数一类 \(t\) 由自然商映射 \(q:P_1=R\to k\) 表示，提升可取底层均为 \(1_R\) 的 \(t_i:P_{i+1}\to P_i\)。它们使次数降低一，而不是原消解中的零次恒等映射。上述乘积公式将 \(t^2\) 表示为 \(q t_1:P_2\to k\)，它仍是非零商映射；Hom 复形的微分为零，故它也不是余边。因此 \(t^2\ne0\)。若 \(\operatorname{char}k\ne2\)，分次交换律却会要求 \(t^2=-t^2\)，这说明正次数中的障碍也确实存在。\cref{b4:ex:09-dual-number-product}进一步计算这一例子的整个自扩张环。

\ProseParagraphBreak
消解乘积的计算采用 \(\delta f=f\mathrm{d}\) 的次数约定。与\cref{b4:def:hom-complex}的内部 Hom 约定比较时，把 \(P_i,Q_i\) 放在上链次数 \(-i\)，并记 \(c_r=(-1)^{r(r+1)/2}\)。\secref{b4:sec:0802}的重标度将内部 Hom 的余圈 \(c_nb\) 对应于这里的 \(b\)。相应的次数 \(n\) 提升在分量 \(P_{n+i}\to Q_i\) 上须取
\[
 T_i=c_n(-1)^{ni}t_i,\qquad
 \mathrm{d}_QT_i=(-1)^n\cdot T_{i-1}\mathrm{d}_P\quad(i\ge1),
\]
才能满足内部 Hom 的余圈条件。将另一个余圈写成 \(c_ma\)，其复合为 \(c_mc_n(-1)^{mn}at_m\)。由于
\[
 c_{m+n}=(-1)^{mn}c_mc_n,
\]
在次数 \(m+n\) 作回同一重标度后，仍得到 \(at_m\)。两种符号约定因此给出同一个乘积，但须同时调整余圈与全部提升分量。Ext 扩张的基本对应和 Baer 和可参阅 Weibel\cite[Theorem 3.4.3, Definition 3.4.4, Corollary 3.4.5]{Weibel1994HomologicalAlgebra}；高阶扩张的进一步讨论见该书\cite[Vista 3.4.6]{Weibel1994HomologicalAlgebra}。

% !TeX root = ../../Book4.tex
% 纲要标题：### 10.2 截短消解与同调维数估计
% 纲要位置：计划纲要.md，第 2719 行。

\section{截短消解与同调维数估计}
\label{b4:sec:1002}

高次导出群是否消失，可以转化为消解末端的合冲模是否投射。这个判据既判断消解能否终止，也说明局部化以后投射维数需要怎样核对。

% 纲要内容：
%
% * 截短消解及合冲模
% * 短正合列中的同调维数不等式
% * 第三卷局部投射维数所需的最低平移版本已在第8.5节证明；本节作一般化整理
% * 局部投射维数的讨论作为交换代数应用，可在学习第三卷后回读；不参与前两小节的证明
%

\subsection{截短消解及合冲模}
\label{b4:subsec:1002:01}

第八章已经证明一步维数平移：短正合列中的投射或内射对象使长正合列的相邻高次项同构。本节用它确定消解何时真正能够终止，而不再重复连接同态的构造。

\begin{proposition}\label{b4:prop:truncation-dimension}
设 \(R\) 为含幺环，\(M\) 为幺性左 \(R\) 模，\(d\ge1\) 为整数。给定左 \(R\) 模正合列
\[
 0\longrightarrow K_d\longrightarrow P_{d-1}\longrightarrow\cdots
 \longrightarrow P_0\longrightarrow M\longrightarrow0,
\]
其中各 \(P_i\) 投射。则 \(\operatorname{pd}_RM\le d\) 当且仅当 \(K_d\) 投射。若 \(0\to M\to I^0\to\cdots\to I^{d-1}\to L^d\to0\) 正合且各 \(I^i\) 内射，则 \(\operatorname{id}_RM\le d\) 当且仅当 \(L^d\) 内射。若第一列的 \(P_i\) 仅要求平坦，则 \(\operatorname{fd}_RM\le d\) 当且仅当 \(K_d\) 平坦。
\end{proposition}
\begin{proof}
若末端对象具有所需性质，给定列本身就是相应长度的消解。反向分别利用
\[
 \operatorname{Ext}_R^1(K_d,N)\simeq\operatorname{Ext}_R^{d+1}(M,N),\qquad
 \operatorname{Ext}_R^1(N,L^d)\simeq\operatorname{Ext}_R^{d+1}(N,M),
\]
以及 \(\operatorname{Tor}_1^R(T,K_d)\simeq\operatorname{Tor}_{d+1}^R(T,M)\)。投射和内射的式子是\cref{b4:prop:iterated-dimension-shift}；平坦版本对各短正合列用 Tor 长正合列，并用平坦项高次 Tor 为零。右侧由\cref{b4:thm:dimension-vanishing}消失，左侧的次数一判据给出所需性质。
\end{proof}

因此某个 \(d\) 步合冲模投射时，任何投射消解的 \(d\) 步合冲模都投射，但这些合冲模不必同构。给一个消解额外加入自由直和项，就会改变实际的核；不变的是上述消失条件及其表达的维数。

\subsection{短正合列中的同调维数不等式}
\label{b4:subsec:1002:02}

\begin{proposition}\label{b4:prop:dimension-short-exact}
设 \(R\) 为含幺环，\(0\to A\to B\to C\to0\) 为幺性左 \(R\) 模短正合列。以下所有维数均在 \(R\) 上计算。投射维数满足
\[
\begin{aligned}
 \operatorname{pd}B&\le\max\{\operatorname{pd}A,\operatorname{pd}C\},\\
 \operatorname{pd}A&\le\max\{\operatorname{pd}B,\operatorname{pd}C-1\},\\
 \operatorname{pd}C&\le\max\{\operatorname{pd}B,\operatorname{pd}A+1\}.
\end{aligned}
\]
平坦维数满足相同的三式。内射维数满足
\[
\begin{aligned}
 \operatorname{id}B&\le\max\{\operatorname{id}A,\operatorname{id}C\},\\
 \operatorname{id}A&\le\max\{\operatorname{id}B,\operatorname{id}C+1\},\\
 \operatorname{id}C&\le\max\{\operatorname{id}B,\operatorname{id}A-1\}.
\end{aligned}
\]
这里非零模的维数非负，零模的维数为 \(-\infty\)，并按通常扩展约定处理 \(\pm\infty\) 与有限整数的加减。
\end{proposition}
\begin{proof}
固定左模 \(N\)。Ext 的第一变量长正合列包含
\[
 \operatorname{Ext}^i(C,N)\to\operatorname{Ext}^i(B,N)\to\operatorname{Ext}^i(A,N)
 \to\operatorname{Ext}^{i+1}(C,N).
\]
第一式由两端的消失得到中间的消失；第二式用 \(\operatorname{Ext}^i(B,N)\) 及 \(\operatorname{Ext}^{i+1}(C,N)\)；第三式把次数改成 \(i-1,i\)，使用 \(\operatorname{Ext}^{i-1}(A,N)\) 及 \(\operatorname{Ext}^i(B,N)\)。对每个大于相应上界的正整数次数，这些群均为零，再用\cref{b4:thm:dimension-vanishing}。若上界为负，则只能来自零对象或 \(B=0\) 的平凡短列，结论直接成立。

平坦情形固定右模 \(T\)，用 \(\operatorname{Tor}_i(T,-)\) 的长正合列中三种相邻位置得到相同上界。内射情形改用第二变量长正合列
\(\operatorname{Ext}^i(N,A)\to\operatorname{Ext}^i(N,B)\to\operatorname{Ext}^i(N,C)\to\operatorname{Ext}^{i+1}(N,A)\)，逐一取相邻项，得到所写三式。
\end{proof}

这些不等式也给出精确的一步下降。例如 \(0\to K\to P\to M\to0\) 中 \(P\) 投射，若 \(\operatorname{pd}M=d\ge1\) 有限，则 \(\operatorname{pd}K=d-1\)：一式给出不超过 \(d-1\)，另一式排除更小值。若 \(d=0\)，短列分裂，\(K\) 投射或为零；不能强行套用数值 \(d-1=-1\)。

\subsection{交换代数应用：局部投射维数}
\label{b4:subsec:1002:03}

前两小节的截短判据与维数不等式已经由第八章的长正合列和维数平移、\secref{b4:sec:1001}的消失判据证明。本小节把这些一般结果用于交换 Noether 局部环；所调用的第三卷极小消解理论只参与这项应用，不构成前两小节的前提。

第三卷命题17.1.2已经在交换 Noether 局部环上构造有限生成模的极小自由消解，并证明其微分矩阵的元素都在极大理想中。定理17.2.3据此建立局部投射维数判据。这里用本节的一般理论解释其中“某一步消失就能截短”的依据。

\ProseParagraphBreak
具体地，设 \((R,\mathfrak m,k)\) 为交换 Noether 局部环，\(M\ne0\) 为有限生成的 \(R\)-模，\(F_\bullet\to M\) 为极小自由消解。第三卷命题17.1.3的极小消解与 Betti 数\footnote{贝蒂（Enrico Betti，1823-1892），意大利数学家，研究代数与拓扑；以其名字命名的数记录各维同调的秩。}结论给出
\[
 \operatorname{Tor}_i^R(k,M)\simeq k\otimes_R F_i,\qquad
 \operatorname{Ext}_R^i(M,k)\simeq\operatorname{Hom}_R(F_i,k).
\]
理由是两个复形的微分都变为零。有限性使 Nakayama 引理适用，所以 \(k\otimes F_i=0\) 就是 \(F_i=0\)，此后极小消解终止。结合本节的截短判据，投射维数等于这些 Tor 非零次数的上确界。这解释了为什么在这一特定环境中只需测试剩余域，而一般环上的判据必须量化所有模。

\ProseParagraphBreak
例如在 \(R=k[x,y]_{(x,y)}\) 上，剩余域的 Koszul 消解为
\[
 0\longrightarrow R\xrightarrow{\binom{-y}{x}}R^2
 \xrightarrow{(x\ \ y)}R\longrightarrow k\longrightarrow0.
\]
它极小，故 \(\operatorname{Tor}_2^R(k,k)=k\ne0\)，同时长度二给出高次消失，从而 \(\operatorname{pd}_Rk=2\)。相反，\cref{b4:exm:infinite-dimension-dual-numbers}的每一步都保留一个自由项，有限生成并不意味着有限投射维数。这里不重证第三卷的极小消解存在性或 Auslander--Buchsbaum 公式\footnote{奥斯兰德（Maurice Auslander，1926-1994），美国数学家，研究同调代数与表示论；布赫斯鲍姆（David Alvin Buchsbaum，1929-2021），美国数学家，研究同调代数及交换代数中的自由消解。}；它们分别依赖局部有限生成理论和深度理论。

% !TeX root = ../../Book4.tex
% 纲要标题：### 10.3 万有系数定理
% 纲要标签：b4:sec:0903
% 纲要位置：计划纲要.md，第 2697 行。

\section{万有系数定理}
\label{b4:sec:0903}

改变系数可能使原来的单射失效，因此同调与系数张量之间会多出 Tor 项。主理想整环上圈与边界自由，使这个差额形成自然的短正合列。

% 纲要内容：
%
% * 主理想整环上的万有系数短正合列
% * 自然短正合列与非典范分裂
% * 系数变换的计算
%

\subsection{主理想整环上的万有系数短正合列}
\label{b4:subsec:0903:01}

主理想整环上，自由模的子模仍自由，即使自由模的秩无限也成立，见\cref{b4:thm:pid-free-submodule}。因此自由链复形的圈和边界都是自由模。这个事实既使 Tor 项可由长度一的消解算出，也保证下面的短正合列能够分裂。

\begin{theorem}[同调的万有系数定理]\label{b4:thm:uct-homology}
设 \(R\) 为主理想整环，\(C\) 为自由 \(R\)-模链复形，\(A\) 为 \(R\)-模。对每个整数 \(n\)，有短正合列
\[
 0\longrightarrow H_n(C)\otimes_R A
 \xrightarrow{\kappa}H_n(C\otimes_R A)
 \xrightarrow{\rho}\operatorname{Tor}_1^R(H_{n-1}(C),A)
 \longrightarrow0.
\]
此列对自由链复形之间的链映射 \(C\to C'\) 和系数模同态 \(A\to A'\) 都协变自然。它分裂，但一般没有关于 \(C\) 自然的分裂。
\end{theorem}
\begin{proof}
令 \(Z_n=\ker \mathrm{d}_n\)、\(B_n=\operatorname{im}\mathrm{d}_{n+1}\)。由于 \(B_{n-1}\) 自由，序列 \(0\to Z_n\to C_n\to B_{n-1}\to0\) 分裂。张量后仍是短正合列，并拼成复形的短正合列，其中外侧复形的微分为零。其同调长正合列在所需位置给出
\[
 B_n\otimes A\longrightarrow Z_n\otimes A\longrightarrow H_n(C\otimes A)
 \longrightarrow B_{n-1}\otimes A\longrightarrow Z_{n-1}\otimes A.
\]
末箭头是包含 \(B_{n-1}\hookrightarrow Z_{n-1}\) 张量 \(A\)，这是由连接同态对提升的微分直接算出的。前一箭头的余核是 \(H_n(C)\otimes A\)；后一箭头的核由自由消解 \(0\to B_{n-1}\to Z_{n-1}\to H_{n-1}(C)\to0\) 识别为所写的 Tor 群。所有映射在选分裂之前已经定义，所以序列自然。

固定次数 \(n\)，选取投影 \(p_n:C_n\to Z_n\)，并记商映射为 \(q_n:Z_n\to H_n(C)\)。在 \(C\otimes A\) 的 \(n\) 次圈上定义
\[
 [x]\longmapsto(q_n\otimes1_A)(p_n\otimes1_A)(x).
\]
若 \(x\) 改变一个边界 \((\mathrm{d}_{n+1}\otimes1_A)(y)\)，由于 \(p_n\) 在 \(Z_n\) 上是恒等映射，而 \(\mathrm{d}_{n+1}\) 的像在 \(B_n\subseteq Z_n\) 中，这个差经 \(q_n\otimes1_A\) 映为零。因此公式在同调上良定义；它与 \(\kappa\) 的复合是恒等映射，故给出分裂。

这里的投影只在固定次数上使用。即使逐次选取 \(p_n\)，它们一般也不构成从 \(C\) 到零微分复形 \(Z\) 的链映射：\(p_{n-1}\mathrm{d}_n=\mathrm{d}_n\)，而链映射条件要求这个复合为零。上述分裂依赖补模的选择；下一小节给出无法随链复形自然选取分裂的例子。
\end{proof}

\begin{theorem}[上同调的万有系数定理]\label{b4:thm:uct-cohomology}
设 \(R\) 为主理想整环，\(C\) 为自由 \(R\)-模链复形，\(A\) 为 \(R\)-模。令上链复形 \(\operatorname{Hom}_R(C,A)^n=\operatorname{Hom}_R(C_n,A)\)，微分为 \(f\mapsto f\circ\mathrm{d}_{C,n+1}\)。则对每个整数 \(n\)，有短正合列
\[
 0\longrightarrow\operatorname{Ext}_R^1(H_{n-1}(C),A)
 \longrightarrow H^n(\operatorname{Hom}_R(C,A))
 \xrightarrow{r}\operatorname{Hom}_R(H_n(C),A)\longrightarrow0.
\]
此列对自由链复形之间的链映射 \(C\to C'\) 反变自然，对系数模同态 \(A\to A'\) 协变自然；前者在中间项诱导从 \(H^n(\operatorname{Hom}_R(C',A))\) 到 \(H^n(\operatorname{Hom}_R(C,A))\) 的映射。此列分裂，但一般没有关于 \(C\) 自然的分裂。
\end{theorem}
\begin{proof}
余圈 \(f:C_n\to A\) 零化 \(B_n\)，所以限制到 \(Z_n\) 后经 \(H_n(C)\) 分解；余边的限制为零，这定义了 \(r\)。由于 \(Z_n\) 在 \(C_n\) 中为直和因子，任意 \(H_n(C)\to A\) 都能先拉回 \(Z_n\) 再延拓至 \(C_n\)，故 \(r\) 满射。

\(r([f])=0\) 恰好表示 \(f|_{Z_n}=0\)，因而 \(f=t \mathrm{d}_n\)，其中 \(t:B_{n-1}\to A\)。这样的 \(f\) 是余边，当且仅当 \(t\) 能延拓至 \(C_{n-1}\)。因为 \(Z_{n-1}\) 也是直和因子，这又等价于 \(t\) 能延拓至 \(Z_{n-1}\)。于是
\[
 \ker r=\operatorname{Hom}_R(B_{n-1},A)/
 \operatorname{res}\operatorname{Hom}_R(Z_{n-1},A)
 =\operatorname{Ext}_R^1(H_{n-1}(C),A).
\]
链映射保持圈与边界，系数同态则与上述限制、分解及取商相容，故这个核的识别给出所述自然性。

选定投影 \(p_n:C_n\to Z_n\)，并记 \(q_n:Z_n\to H_n(C)\) 为商映射。对 \(\phi:H_n(C)\to A\)，取 \(n\) 次余链 \(\phi q_np_n\)。由于 \(p_n\) 在 \(Z_n\) 上是恒等映射，且 \(q_n\) 零化 \(B_n=\mathrm{d}_{n+1}(C_{n+1})\)，有
\[
 (\phi q_np_n)\mathrm{d}_{n+1}=0.
\]
所以它是余圈；限制到 \(Z_n\) 后为 \(\phi q_n\)，从而 \(\phi\mapsto[\phi q_np_n]\) 是 \(r\) 的右逆。这里构造的是固定次数上的余圈代表，无须把投影 \(p_n\) 本身视为链映射。
\end{proof}

上述 Hom 计算与\cref{b4:def:hom-complex}的内部 Hom 仅差每度的符号重标识，所得上同调相同。两个版本可与 Weibel\cite[Universal Coefficient Theorem for Homology 3.6.2, Universal Coefficient Theorem for Cohomology 3.6.5]{Weibel1994HomologicalAlgebra}对照；该书把同调版本先写在自由交换群上，本节使用同一证明得到主理想整环版本。

\subsection{自然短正合列与非典范分裂}
\label{b4:subsec:0903:02}

“自然短正合列”是指链映射 \(C\to C'\) 及系数同态 \(A\to A'\) 诱导整个正合列的交换图：同调版本对两变量都协变，上同调版本对 \(C\) 反变、对 \(A\) 协变。它不意味着左右两项在中间项中都已选定。自然的是第一项的像及第二项的商，而直接把中间项写成直和还要作额外选择。

\begin{example}\label{b4:exm:uct-no-natural-splitting}
下面固定系数 \(A=\mathbb F_2\)，说明分裂不能对链复形变量 \(C\) 自然。在 \(\mathbb Z\) 上取 \(C_1=\mathbb Za\oplus\mathbb Zb\)、\(C_0=\mathbb Zc\)，\(\mathrm{d}(a)=2c\)、\(\mathrm{d}(b)=0\)，其余项为零。链自同构
\[
 f_1(a)=a+b,\quad f_1(b)=b,\quad f_0(c)=c
\]
在 \(H_1(C)=\mathbb Zb\) 及 \(H_0(C)=\mathbb Z/2\) 上都诱导恒等映射，而在
\(H_1(C\otimes A)=A\bar a\oplus A\bar b\) 上是非平凡剪切。同调万有系数列左项是 \(A\bar b\)，商由 \(\bar a\) 表示。任何商的截面都把 \(1\) 送到 \(\bar a+\lambda\bar b\)，但 \(f\) 把它改成 \(\bar a+(\lambda+1)\bar b\)。所以没有与全部链映射相容的截面。

\ProseParagraphBreak
上同调版本也由同一例子说明。在 \(H^1(\operatorname{Hom}(C,A))\) 中，左端 Ext 项由 \(a^*\) 生成，右端 Hom 项由 \(b^*\) 表示。\(f^*(a^*)=a^*\)、\(f^*(b^*)=a^*+b^*\)，再一次排除了自然截面。
\end{example}

这个反例没有排除固定 \(C\) 后对系数变量的自然性。事实上，一旦固定各次数上的投影 \(p_n:C_n\to Z_n\)，前述两种分裂公式都与系数同态 \(A\to A'\) 相容。

\ProseParagraphBreak
如果 \(A\) 平坦，同调公式的 Tor 项消失，\(\kappa\) 就是自然同构；如果 \(H_{n-1}(C)\) 自由，两个版本的附加项都消失。这两种简化来自明确的消失条件，不需要选择分裂。

\subsection{系数变换的计算}
\label{b4:subsec:0903:03}

设 \(C\) 是有限秩自由整数链复形。若已由 Smith 标准形\footnote{史密斯（Henry John Stephen Smith，1826-1883），爱尔兰数学家，研究二次型、初等因子与数论。}算得
\[
 H_n(C)\simeq\mathbb Z^{r_n}\oplus\bigoplus_j\mathbb Z/d_{n,j},
 \qquad d_{n,j}>1,
\]
则改变系数不必重新化简整个张量复形。对 \(A=\mathbb Z/m\)，令 \(g_{n,j}=\gcd(d_{n,j},m)\)。长度一的循环模消解给出
\[
 (\mathbb Z/d)\otimes\mathbb Z/m\simeq\mathbb Z/\gcd(d,m),\qquad
 \operatorname{Tor}_1^{\mathbb Z}(\mathbb Z/d,\mathbb Z/m)
 \simeq\mathbb Z/\gcd(d,m).
\]
因此同调的同构类型是
\[
 H_n(C\otimes\mathbb Z/m)\simeq
 (\mathbb Z/m)^{r_n}\oplus\bigoplus_j\mathbb Z/g_{n,j}
 \oplus\bigoplus_j\mathbb Z/g_{n-1,j}.
\]
这个式子确定模的同构类型，并未指定内部直和因子。

\begin{example}\label{b4:exm:uct-integer-computation}
取 \(C_2=\mathbb Z\)、\(C_1=\mathbb Z^2\)、\(C_0=\mathbb Z\)，
\[
 \mathrm{d}_2(t)=(0,6t),\qquad \mathrm{d}_1(x,y)=4x.
\]
有 \(H_2(C)=0\)、\(H_1(C)=\mathbb Z/6\)、\(H_0(C)=\mathbb Z/4\)。取 \(A=\mathbb Z/2\)，两个微分都变成零，故新的同调依次为 \(A,A^2,A\)（次数 \(2,1,0\)）。在次数 \(1\)，其中一份 \(A\) 来自 \(H_1(C)\otimes A\)，另一份来自 \(\operatorname{Tor}_1(H_0(C),A)\)；次数 \(2\) 完全来自 \(\operatorname{Tor}_1(H_1(C),A)\)。

\ProseParagraphBreak
取上同调的整数系数，则 \(H^0(\operatorname{Hom}(C,\mathbb Z))=0\)、\(H^1=\mathbb Z/4\)、\(H^2=\mathbb Z/6\)。这些群都来自相邻次数的 Ext 项；直接对两个微分作转置也得到同样结果。因而“对偶以后只需对同调取 Hom”会漏去本例全部的非零上同调。
\end{example}

% !TeX root = ../../Book4.tex
% 纲要标题：### 10.4 Künneth 公式
% 纲要标签：b4:sec:0904
% 纲要位置：计划纲要.md，第 2703 行。

\section{Künneth 公式}
\label{b4:sec:0904}

张量两个复形的同调，由各自同调的张量积及一次 Tor 共同控制。域上 Tor 消失，而主理想整环上即使短正合列分裂，分裂也通常没有自然选择。

% 纲要内容：
%
% * 域和主理想整环的基本版本
% * Tor 项及成立假设
% * 与第三卷 Koszul 计算的联系
%
% ---
%

\subsection{域和主理想整环的基本版本}
\label{b4:subsec:0904:01}

\begin{theorem}[Künneth 公式]\label{b4:thm:kunneth-pid}
设 \(R\) 为主理想整环，\(C,D\) 为下有界的自由 \(R\)-模链复形。对每个整数 \(n\)，有对 \(C,D\) 的链映射自然的短正合列
\[
\begin{split}
0\longrightarrow&\ \bigoplus_{p+q=n}H_p(C)\otimes_R H_q(D)
 \xrightarrow{\kappa_n}H_n(C\otimes_R D)\\
 \longrightarrow&\ \bigoplus_{p+q=n-1}\operatorname{Tor}_1^R(H_p(C),H_q(D))
 \longrightarrow0.
\end{split}
\]
此列非典范地分裂。特别地，若 \(R=k\) 为域，则 \(\kappa_n\) 是自然同构。
\end{theorem}
\begin{proof}
把 \(Z_p=Z_p(C)\) 放在次数 \(p\)，组成零微分复形 \(Z\)；把 \(B_{p-1}=B_{p-1}(C)\) 放在次数 \(p\)，组成零微分复形 \(B'\)。短正合列 \(0\to Z\to C\to B'\to0\) 逐项分裂，因此与 \(D\) 张量后仍是短正合列。\(Z_p,B_p\) 都自由，由\cref{b4:prop:flat-coefficients-homology}，两侧复形的同调分别是
\[
 \bigoplus_{p+q=n}Z_p\otimes H_q(D),\qquad
 \bigoplus_{p+q=n}B_{p-1}\otimes H_q(D).
\]
在同调长正合列中，连接同态在每一项上由包含 \(B_{p-1}\hookrightarrow Z_{p-1}\) 张量 \(H_q(D)\) 给出。具体地，把边界元素提升到 \(C_p\)，再作总微分，来自 \(D\) 的一项对圈代表为零，剩下的就是该包含。这些包含张量映射的余核与核，由自由消解
\(0\to B_p\to Z_p\to H_p(C)\to0\) 分别识别为张量积及 \(\operatorname{Tor}_1\)，从而得到所写自然正合列。

还须证明分裂。记 \(i_p:B_p\hookrightarrow Z_p\) 为包含。因为 \(B_{p-1}\) 自由，可以为满射 \(\mathrm d_p:C_p\to B_{p-1}\) 选取截面 \(s_{p-1}:B_{p-1}\to C_p\)。于是
\[
 Z_p\oplus B_{p-1}\xrightarrow{\ \sim\ }C_p,
 \qquad (z,b)\longmapsto z+s_{p-1}(b).
\]
在这些坐标下，\(\mathrm d_p(z,b)=(i_{p-1}(b),0)\in Z_{p-1}\oplus B_{p-2}\)。因此 \(C\) 同构于两项复形的直和
\[
 C\simeq\bigoplus_p E(p),\qquad
 E(p):\quad B_p\xrightarrow{i_p}Z_p
 \quad(\text{次数 }p+1,p),
\]
其中 \(E(p)\) 的其余项为零。对 \(D\) 同样选取截面，得到
\[
 D\simeq\bigoplus_q F(q),\qquad
 F(q):\quad B_q(D)\hookrightarrow Z_q(D)
 \quad(\text{次数 }q+1,q),
\]
其中 \(F(q)\) 的其余项为零，且同调仅可能在次数 \(q\) 非零，其值为 \(H_q(D)\)。于是张量复形分解为 \(E(p)\otimes_R F(q)\) 的直和；下有界性保证每个总次数只涉及有限多个非零的 \((p,q)\)。

把 \(H_q(D)\) 看成只在次数 \(q\) 非零的复形。商映射 \(Z_q(D)\to H_q(D)\) 给出拟同构 \(F(q)\to H_q(D)\)。因为 \(E(p)\) 有界且自由，\cref{b4:prop:bounded-flat-tensor}说明
\[
 E(p)\otimes_R F(q)\longrightarrow E(p)\otimes_R H_q(D)
\]
也是拟同构。右侧是 \(B_p\otimes_R H_q(D)\to Z_p\otimes_R H_q(D)\)，放在次数 \(p+q+1,p+q\)。由 \(H_p(C)\) 的长度一自由消解，它在这两个次数的同调分别为 \(\operatorname{Tor}_1^R(H_p(C),H_q(D))\) 与 \(H_p(C)\otimes_R H_q(D)\)。因此在总次数 \(n\)，张量积部分来自 \(p+q=n\)，Tor 部分来自 \(p+q=n-1\)，得到直和分解。圈的张量代表元表明，前一部分恰是 \(\kappa_n\) 的像，后一部分通过自然正合列的商映射同构到其右端，因而给出一个分裂。

自然性来自前半段的圈模、边界模及同调长正合列：链映射保持这些结构，所以 \(\kappa_n\) 和右端商映射都是自然的。后半段所选截面却一般不与链映射相容，故所得直和分解没有自然选择。这与万有系数定理中自然短正合列和所选分裂的区别相同。

域上所有模平坦，Tor 项为零，故最后的结论成立。
\end{proof}

更一般的 Künneth 公式见 Weibel\cite[Theorem 3.6.3]{Weibel1994HomologicalAlgebra}：只要求一个复形的各项及边界模平坦，另一个复形可以任意，也不要求两者下有界。本节采用下有界自由复形这一较强的充分条件，使每条总次数对角线有限，并可直接沿用\cref{b4:prop:bounded-flat-tensor}的总化论证；下有界性并非 Künneth 公式本身的必要条件。

\subsection{Tor 项及成立假设}
\label{b4:subsec:0904:02}

\begin{example}\label{b4:exm:kunneth-two-cyclic}
设 \(m,n\) 为正整数，在 \(\mathbb Z\) 上分别取两项自由消解 \(C=(\mathbb Z\xrightarrow{m}\mathbb Z)\)、\(D=(\mathbb Z\xrightarrow{n}\mathbb Z)\)，次数均为 \(1,0\)。张量总复形是
\[
 0\longrightarrow\mathbb Z\xrightarrow{\binom{-n}{m}}
 \mathbb Z^2\xrightarrow{(m\ \ n)}\mathbb Z\longrightarrow0.
\]
令 \(g=\gcd(m,n)\)。第一个映射单射；第二个映射的像为 \(g\mathbb Z\)，核由 \((n/g,-m/g)\) 生成，而第一个映射的像是该核的 \(g\) 倍。因此
\[
 H_0(C\otimes D)=\mathbb Z/g,\qquad H_1(C\otimes D)=\mathbb Z/g,
 \qquad H_2(C\otimes D)=0.
\]
第二个非零群是 Tor 项；只取两个消解原有同调的张量积，会漏掉它。
\end{example}

定理中的自由性假设不能直接删去。仍取主理想整环 \(R=\mathbb Z\)，让 \(C,D\) 都是集中在次数零的 \(\mathbb Z/2\)。此时普通张量复形仅在次数零非零，但 \(\operatorname{Tor}_1^{\mathbb Z}(\mathbb Z/2,\mathbb Z/2)=\mathbb Z/2\)。若无条件套用上面的短正合列，在次数一会错误地要求零模满射到 \(\mathbb Z/2\)。这里应先取平坦消解再张量；普通张量复形本身没有包含该高阶信息。

\ProseParagraphBreak
主定理中的下有界性保证每条总次数对角线只有有限个非零项，因而上述有限对角总化工具适用。在一般环上，若要用无界复形替换张量积中的某个因子，还须检验它是否保持拟同构；逐项自由本身不能保证这一点，见\cref{b4:exm:unbounded-flat-tensor}。该反例的底环不是主理想整环，不能据此认定主定理中的下有界性是必要条件。

\subsection{与第三卷 Koszul 计算的联系}
\label{b4:subsec:0904:03}

第三卷第十五章的 Koszul 复形，可以看成两项复形的张量积。对交换环 \(R\) 中元素 \(x\)，写
\[
 K(x):\quad 0\longrightarrow R\xrightarrow{x}R\longrightarrow0
 \qquad(\text{次数 }1,0).
\]
多个元素的 Koszul 复形由 \(K(x_1)\otimes_R\cdots\otimes_R K(x_r)\) 给出。各因子次数一生成元的交换引入负号，这与外代数描述中的符号相同。第三卷命题15.2.2已经证明该张量识别，第15.3节证明正则序列时的正合性；这里说明换系数后如何读出新的同调。

\begin{proposition}\label{b4:prop:koszul-adjoin-element}
设 \(R\) 为交换含幺环，\(C\) 为 \(R\) 模链复形，\(x\in R\)。记 \(K(x)\) 为仅在次数 \(1,0\) 非零的链复形 \(R\xrightarrow{x}R\)，微分为乘 \(x\)。对每个整数 \(n\)，有自然短正合列
\[
 0\longrightarrow H_n(C)/xH_n(C)
 \longrightarrow H_n(C\otimes_R K(x))
 \longrightarrow\{\,z\in H_{n-1}(C):xz=0\,\}
 \longrightarrow0.
\]
这里不要求 \(R\) 为主理想整环，也不保证分裂。
\end{proposition}
\begin{proof}
张量复形在次数 \(n\) 是 \(C_n\oplus C_{n-1}\)，其微分为
\(\mathrm{d}(a,b)=(\mathrm{d}_Ca+(-1)^{n-1}xb,\mathrm{d}_Cb)\)。取包含 \(a\mapsto(a,0)\) 及投影 \(q_n(a,b)=(-1)^{n-1}b\)，便得到短正合列
\(0\to C\to C\otimes K(x)\xrightarrow{q}C[1]\to0\)，其中链记号下 \(C[1]_n=C_{n-1}\)、\(\mathrm{d}_{C[1]}=-\mathrm{d}_C\)。等式 \(q_{n-1}\mathrm{d}=-\mathrm{d}_Cq_n\) 直接保证投影是链映射。对于 \(C_{n-1}\) 中的圈 \(z\)，提升为 \((0,(-1)^{n-1}z)\)，微分是 \((xz,0)\)，故连接同态就是乘 \(x\)。因此同调长正合列的相关部分是
\[
 H_n(C)\xrightarrow{x}H_n(C)\longrightarrow H_n(C\otimes K(x))
 \longrightarrow H_{n-1}(C)\xrightarrow{x}H_{n-1}(C).
\]
取中间部分的核与余核即得结论。
\end{proof}

例如在 \(R=k[x]\) 上取两个相同元素 \(x,x\)。第一个复形 \(K(x)\) 只有 \(H_0=k\)，而 \(x\) 在该同调上作用为零，故 \(K(x,x)\) 的 \(H_0,H_1\) 都是 \(k\)。次数一的具体圈是 \(e_1-e_2\)，其倍数 \(x(e_1-e_2)\) 是次数二的边界（依生成元次序差一个负号）。这说明第二个 \(x\) 不是 \(R/(x)\) 上的非零因子，故重复元素不能形成正则序列。


\begin{chaptersummary}
扩张把 Ext 的余圈解释为正合列，投射消解的截短与推出使两种描述可以互相转换。Baer 和来自对角与加法，Yoneda 乘积来自拼接；它们分别对应余圈的相加和提升后的复合。

\ProseParagraphBreak
张量总复形中的符号保证微分平方为零，平坦性和有界性则承担保持同调的不同任务。在主理想整环上，自由子模的结构使万有系数与 Künneth 公式只需一次 Ext 或 Tor 修正。自然的是短正合列及其映射，分裂通常依赖选择。循环模的两项消解和 Koszul 复形都提供了能够直接核验这些差别的计算。
\end{chaptersummary}

\BookExercises

\begin{exercise}\label{b4:ex:09-extensions-parameters}
设 \(p\) 为素数，\(t\in\mathbb F_p\)。用生成元 \(x,y\) 和关系 \(px=0\)、\(py=tx\) 定义交换群 \(E_t\)。定义 \(i:\mathbb Z/p\mathbb Z\to E_t\) 为 \(i(\bar1)=x\)，定义 \(q:E_t\to\mathbb Z/p\mathbb Z\) 为 \(q(y)=\bar1\)、\(q(x)=\bar0\)。证明 \(0\to\mathbb Z/p\mathbb Z\xrightarrow{i}E_t\xrightarrow{q}\mathbb Z/p\mathbb Z\to0\) 是扩张，求各中间群的同构类型，并说明 Baer 和怎样作用于 \(t\)。
\end{exercise}
\begin{solution}
关系先把 \(b\) 降至 \(0\le b<p\)，再把 \(a\) 降至 \(0\le a<p\)，故每个元素都可写成 \(a x+b y\)。取 \(t\) 的整数代表 \(\widetilde t\)，关系子格由矩阵
\[
 M=\begin{pmatrix}p&-\widetilde t\\0&p\end{pmatrix}
\]
的两列生成。第一卷第 9.4 节的子格指数命题说明，非奇异整数方阵生成的子格在 \(\mathbb Z^2\) 中的指数等于行列式的绝对值；这里 \(|E_t|=|\det M|=p^2\)。上述 \(p^2\) 个代表因而互不相同，表示也唯一。因此 \(x\) 的阶为 \(p\)，核恰为它生成的群。当 \(t=0\) 时 \(E_t=(\mathbb Z/p\mathbb Z)^2\)；当 \(t\ne0\) 时 \(py=tx\ne0\)，所以 \(y\) 的阶为 \(p^2\)，且生成全群。取自由满射 \(\mathbb Z\to\mathbb Z/p\mathbb Z\)，将 \(1\) 提升到 \(y\)，在核 \(p\mathbb Z\) 上对应的映射为 \(p\mapsto t\)。故扩张类为 \(t\)，Baer 和把参数相加。中间群同构并不能消去两端指定映射的区别。
\end{solution}

\begin{exercise}\label{b4:ex:09-pushout-scalar}
设 \(R\) 为交换环，\(a\in R\) 为非零因子，\(N\) 为 \(R\)-模。证明在 \(\operatorname{Ext}_R^1(R/(a),N)\simeq N/aN\) 的识别下，沿乘 \(r:N\to N\) 推出扩张对应乘 \(r\)。扩张类 \(\bar n\) 在何时分裂？
\end{exercise}
\begin{solution}
自由消解为 \(0\to R\xrightarrow{a}R\to R/(a)\to0\)。类 \(\bar n\) 由把关系 \(a\) 送到 \(n\) 的映射表示；推出后这个映射变为把 \(a\) 送到 \(rn\)，故对应 \(\overline{rn}\)。由\cref{b4:thm:ext1-extensions}中分裂扩张对应零元的结论，原扩张分裂当且仅当 \(n\in aN\)。非零因子条件保证所写两项列确实是自由消解。
\end{solution}

\begin{exercise}\label{b4:ex:09-higher-not-isomorphism}
在域 \(k\) 上，比较二阶扩张
\(0\to k\to k^2\to k^2\to k\to0\)，其箭头依次为 \(a\mapsto(a,0)\)、\((x,y)\mapsto(y,0)\)、\((u,v)\mapsto v\)，与 \(0\to k\xrightarrow{1}k\xrightarrow{0}k\xrightarrow{1}k\to0\)。给出扩张态射，并说明为何不能要求中间项逐项同构。
\end{exercise}
\begin{solution}
前一列各位置的核与像都是所写坐标轴，所以正合。在两个中间项分别使用 \((x,y)\mapsto x\)、\((u,v)\mapsto v\)，两端使用恒等映射，得到前一列到后一列的扩张态射；中间方块的两条复合都是零。因此两扩张等价。但中间项的维数分别为二和一，不可能逐项同构。
\end{solution}

\begin{exercise}\label{b4:ex:09-dual-number-product}
设 \(k\) 为域，\(R=k[\varepsilon]/(\varepsilon^2)\)，把 \(k\) 识别为 \(R/(\varepsilon)\)。用周期自由消解计算 \(\operatorname{Ext}_R^*(k,k)\) 的 Yoneda 乘积，证明它是一个次数一元生成的多项式环。若 \(\operatorname{char}k\ne2\)，这说明了什么？
\end{exercise}
\begin{solution}
消解每项为 \(R\)，微分乘 \(\varepsilon\)。施加 \(\operatorname{Hom}_R(-,k)\) 后微分为零，每次 Ext 均为 \(k\)。次数一生成元 \(t\) 由商映射 \(P_1=R\to k\) 表示。为计算乘积，把它提升为
\[
 t_i:P_{i+1}\longrightarrow P_i,\qquad t_i=1_R\quad(i\ge0).
\]
这些映射把消解次数降低一；它们的底层模映射都是恒等，但不是消解的同次数恒等映射。因两侧微分都乘 \(\varepsilon\)，有 \(\mathrm{d}_i t_i=t_{i-1}\mathrm{d}_{i+1}\)（\(i\ge1\)）；底端与增广复合就是代表 \(t\) 的商映射。按\cref{b4:prop:yoneda-resolution-product}的提升复合法，\(t^n\) 由 \(P_n\to k\) 的商映射代表。它非零，且因 Hom 复形的微分为零而不可能是余边，故生成该次数的一维 Ext 群。因此自扩张环为 \(k[t]\)、\(\deg t=1\)。若特征不为二，则 \(t^2\ne0\) 而分次交换律会要求 \(t^2=-t^2\)。这里讨论的是自扩张环的 Yoneda 乘积；即使底环 \(R\) 交换，一般模的自扩张环也不必分次交换。群上同调等场合的分次交换性还使用了相应乘积的额外结构。
\end{solution}

\begin{exercise}\label{b4:ex:09-tensor-sign}
在特征不为二的域 \(k\) 上，令 \(C=D=(k\xrightarrow{1}k)\)，次数为 \(1,0\)。若在张量复形中漏去 \((-1)^p\)，证明“微分”的平方不为零；使用正确符号时证明总复形正合。
\end{exercise}
\begin{solution}
记次数一生成元为 \(u\in C_1\)、\(v\in D_1\)，其微分为次数零生成元 \(u_0,v_0\)。不带符号的第一次像是
\[
 u_0\otimes v+u\otimes v_0
 \in(C_0\otimes D_1)\oplus(C_1\otimes D_0).
\]
再作用一次，两项都映到 \(u_0\otimes v_0\in C_0\otimes D_0\)，故平方为 \(2u_0\otimes v_0\ne0\)。使用正确符号时，按次数一项 \((C_1\otimes D_0)\oplus(C_0\otimes D_1)\) 的坐标次序，微分为 \(k\xrightarrow{(-1,1)}k^2\xrightarrow{(1,1)}k\)，前箭头单射、后箭头满射，中间核由 \((-1,1)\) 生成，故正合。
\end{solution}

\begin{exercise}\label{b4:ex:09-rational-coefficients}
设整数链复形 \(C\) 满足每个 \(H_n(C)\) 都是挠群。证明 \(C\otimes\mathbb Q\) 正合。这里是否必须要求 \(C_n\) 自由？
\end{exercise}
\begin{solution}
\(\mathbb Q\) 作为整数模平坦，\cref{b4:prop:flat-coefficients-homology}给出 \(H_n(C\otimes\mathbb Q)=H_n(C)\otimes\mathbb Q\)。挠元 \(x\) 若被 \(m\ne0\) 零化，则 \(x\otimes q=mx\otimes(q/m)=0\)，故右侧为零。这个证明不要求原复形逐项自由。
\end{solution}

\begin{exercise}\label{b4:ex:09-uct-free-and-torsion}
设有限秩自由整数复形 \(C\) 在相邻次数满足 \(H_n(C)=\mathbb Z^2\oplus\mathbb Z/6\mathbb Z\)、\(H_{n-1}(C)=\mathbb Z\oplus\mathbb Z/10\mathbb Z\)。求 \(H_n(C\otimes\mathbb Z/4\mathbb Z)\) 的同构类型及 \(H^n(\operatorname{Hom}(C,\mathbb Z))\) 的同构类型。
\end{exercise}
\begin{solution}
同调公式的左项为 \((\mathbb Z/4\mathbb Z)^2\oplus\mathbb Z/2\mathbb Z\)，右项为 \(\mathbb Z/2\mathbb Z\)。单凭这两个端项还不能确定中间群的同构类型；这里可用\cref{b4:thm:uct-homology}，因为整数环是主理想整环且 \(C\) 逐项自由，其万有系数短正合列分裂。因此答案是 \((\mathbb Z/4\mathbb Z)^2\oplus(\mathbb Z/2\mathbb Z)^2\)。上同调公式的 Ext 项为 \(\mathbb Z/10\mathbb Z\)，Hom 项为 \(\mathbb Z^2\)；由\cref{b4:thm:uct-cohomology}的分裂，答案是 \(\mathbb Z^2\oplus\mathbb Z/10\mathbb Z\)。这些是同构类型，不指定自然的内部直和。
\end{solution}

\begin{exercise}\label{b4:ex:09-modp-betti}
设 \(C\) 为有限秩自由整数链复形，\(H_n(C)=\mathbb Z^{r_n}\oplus\bigoplus_j\mathbb Z/d_{n,j}\mathbb Z\)。固定素数 \(p\)，以 \(t_n\) 表示被 \(p\) 整除的 \(d_{n,j}\) 的个数。证明 \(\dim_{\mathbb F_p}H_n(C\otimes\mathbb F_p)=r_n+t_n+t_{n-1}\)。
\end{exercise}
\begin{solution}
对每个循环直和项，张量 \(\mathbb F_p\) 及一次 Tor 都是零或 \(\mathbb F_p\)，其非零条件恰为 \(p\mid d_{n,j}\)。自由部分张量贡献 \(r_n\)，而一次 Tor 为零。将这些维数代入万有系数短正合列，维数可加，得到公式；此步不需要选分裂。
\end{solution}

\begin{exercise}\label{b4:ex:09-kunneth-numbers}
取整数两项复形 \(C=(\mathbb Z\xrightarrow{6}\mathbb Z)\)、\(D=(\mathbb Z\xrightarrow{15}\mathbb Z)\)，次数均为 \(1,0\)。求总复形各次同调，并给出一次同调的圈生成元。
\end{exercise}
\begin{solution}
总微分为 \(\mathbb Z\xrightarrow{(-15,6)}\mathbb Z^2\xrightarrow{(6,15)}\mathbb Z\)。后者的像为 \(3\mathbb Z\)，核由 \((5,-2)\) 生成；前者的像由 \(-3(5,-2)\) 生成。因此 \(H_0=H_1=\mathbb Z/3\mathbb Z\)，\(H_2=0\)，一次类由 \([(5,-2)]\) 生成。这分别对应张量项及一次 Tor 项。
\end{solution}

\begin{exercise}\label{b4:ex:09-kunneth-flat-homology}
设 \(R\) 为主理想整环，\(C,D\) 为下有界自由复形，且每个 \(H_p(C)\) 平坦。证明 \(\kappa_n\) 为自然同构。若只知每个 \(C_p\) 平坦，为什么不能作此结论？
\end{exercise}
\begin{solution}
同调平坦使 Künneth 正合列的每个 \(\operatorname{Tor}_1(H_p(C),H_q(D))\) 为零，因此自然映射 \(\kappa_n\) 是同构。逐项平坦则只是复形各项的性质，商 \(Z_p/B_p\) 可能不平坦；上一题的 \(C_p,D_q\) 都自由，但一次同调全由 Tor 产生，\(\kappa_1\) 的定义域为零而值域非零。
\end{solution}

\begin{exercise}\label{b4:ex:09-euler-tensor}
设 \(C,D\) 为同一域上有限维向量空间的有界链复形。记
\[
\chi(C)=\sum_n(-1)^n\dim H_n(C).
\]
证明 \(\chi(C\otimes D)=\chi(C)\chi(D)\)。
\end{exercise}
\begin{solution}
域上的 Künneth 同构给出 \(\dim H_n(C\otimes D)=\sum_{p+q=n}\dim H_p(C)\dim H_q(D)\)。所有求和有限，乘 \((-1)^n\) 求和后以 \((-1)^{p+q}=(-1)^p(-1)^q\) 分离成两个有限和的乘积，得到等式。有界及有限维条件保证欧拉数作为普通整数有定义。
\end{solution}

\begin{exercise}\label{b4:ex:09-koszul-zero-factor}
设 \(R\) 为交换环，\(x\in R\) 为非零因子。求 \(K(x,0)=K(x)\otimes_R K(0)\) 的全部同调，并解释为什么不需要假设 \(R\) 为主理想整环。
\end{exercise}
\begin{solution}
令 \(C=K(x)\)。\(K(0)\) 只在次数 \(1,0\) 各有一个 \(R\)，微分为零，所以
\[
 (C\otimes_RK(0))_n=C_n\oplus C_{n-1},\qquad
 \mathrm{d}(a,b)=(\mathrm{d}_Ca,\mathrm{d}_Cb).
\]
本卷的链移位 \(C[1]_n=C_{n-1}\) 带微分 \(-\mathrm{d}_C\)。因此
\[
 (a,b)\longmapsto(a,(-1)^{n-1}b)
\]
给出 \(C\otimes_RK(0)\to C\oplus C[1]\) 的复形同构：第二分量在次数 \(n-1\) 的符号是 \((-1)^{n-2}\)，恰好等于对次数 \(n\) 的符号再乘负号。\(x\) 为非零因子使 \(C\) 在次数一的同调为零，在次数零的同调为 \(R/(x)\)，故 \(K(x,0)\) 的 \(H_0,H_1\) 都为 \(R/(x)\)，其他同调为零。这个直接复形分解在任意交换环上成立，不需要 Künneth 的主理想整环假设。
\end{solution}

\begin{exercise}\label{b4:ex:09-extension-zero-product}
设 \(R\) 为主理想整环，\(A,B,C\) 为 \(R\)-模。证明任意 \(\alpha\in\operatorname{Ext}_R^1(B,C)\)、\(\beta\in\operatorname{Ext}_R^1(A,B)\) 的 Yoneda 乘积为零，但两因子可以都非零。
\end{exercise}
\begin{solution}
任意 \(R\)-模有长度至多一的自由消解，所以 \(\operatorname{Ext}_R^2(A,C)=0\)，乘积必为零。取 \(R=\mathbb Z\)、\(A=B=C=\mathbb Z/2\mathbb Z\)，一次 Ext 群为 \(\mathbb Z/2\mathbb Z\)，两个因子都可取其非零元。这里的零指拼接所得二阶扩张在 Yoneda 等价下为零，并不要求其中间箭头为零，也不要求中间项与零扩张逐项同构；尤其不意味着两个短扩张至少有一个分裂。
\end{solution}

\begin{exercise}\label{b4:ex:09-uct-coefficient-map}
设 \(m,r\) 为正整数，\(C\) 为自由整数链复形，\(H_n(C)=0\)、\(H_{n-1}(C)=\mathbb Z/m\mathbb Z\)。用上同调万有系数定理求系数同态 \(\mathbb Z\to\mathbb Z/r\mathbb Z\) 在 \(H^n(\operatorname{Hom}(C,-))\) 上诱导的映射。
\end{exercise}
\begin{solution}
没有 Hom 项，故这些群自然等于 \(\operatorname{Ext}_{\mathbb Z}^1(\mathbb Z/m\mathbb Z,A)=A/mA\)。系数同态诱导的映射是 \(\mathbb Z/m\mathbb Z\to(\mathbb Z/r\mathbb Z)/m(\mathbb Z/r\mathbb Z)\)，把整数类送到同一整数的像。右侧同构于 \(\mathbb Z/\gcd(m,r)\mathbb Z\)，所以该映射是相应的自然取商。这里同构来自自然短正合列中另一端为零，不涉及非典范分裂。
\end{solution}

\begin{exercise}\label{b4:ex:09-uct-splitting-coefficients}
设 \(R\) 为主理想整环，\(C\) 为自由 \(R\)-模链复形，固定整数 \(n\) 及投影 \(r_n:C_n\to Z_n(C)\)。对每个 \(R\)-模 \(A\)，记同调万有系数列的映射为 \(\kappa_A,\rho_A\)，记由 \(r_n\) 构造的左逆为
\[
 \lambda_A:H_n(C\otimes_R A)\longrightarrow H_n(C)\otimes_R A.
\]
证明 \(\lambda_A\) 与所有系数同态 \(A\to A'\) 相容，从而固定 \(C\) 后可选取关于 \(A\) 自然的分裂。若换一个投影得到 \(\lambda'_A\)，证明存在唯一的关于 \(A\) 自然的同态
\[
 \theta_A:\operatorname{Tor}_1^R(H_{n-1}(C),A)
 \longrightarrow H_n(C)\otimes_R A,
 \qquad \lambda'_A-\lambda_A=\theta_A\rho_A.
\]
\end{exercise}
\begin{solution}
\(\lambda_A\) 由 \(r_n\otimes1_A\) 后接圈模到同调模的商映射诱导；这两个映射都与系数同态相容，所以所得左逆关于 \(A\) 自然。\(\rho_A\) 限制为 \(\ker\lambda_A\) 到右端 Tor 群的同构：它的核是 \(\ker\lambda_A\cap\operatorname{im}\kappa_A=0\)；若 \(\rho_A(x)=t\)，则 \(x-\kappa_A\lambda_A(x)\) 属于 \(\ker\lambda_A\) 且仍映到 \(t\)。这些限制同构与系数同态相容，其逆给出关于 \(A\) 自然的截面。

因为两个左逆在 \(\operatorname{im}\kappa_A=\ker\rho_A\) 上相同，\((\lambda'_A-\lambda_A)\kappa_A=0\)，差映射唯一经余核 \(\rho_A\) 分解，得到 \(\theta_A\)。对系数同态 \(g:A\to A'\)，记中间同调映射为 \(h_g=H_n(1_C\otimes g)\)，右端 Tor 映射为 \(T_g=\operatorname{Tor}_1(H_{n-1}(C),g)\)。由已有映射的自然性，
\[
 \begin{aligned}
 \theta_{A'}T_g\rho_A
 &=\theta_{A'}\rho_{A'}h_g
 =(\lambda'_{A'}-\lambda_{A'})h_g\\
 &=(1\otimes g)(\lambda'_A-\lambda_A)
 =(1\otimes g)\theta_A\rho_A.
 \end{aligned}
\]
再用 \(\rho_A\) 的满射性消去最右侧的 \(\rho_A\)，即得 \(\theta_{A'}T_g=(1\otimes g)\theta_A\)，这正是 \(\theta\) 的自然性。这里固定了 \(C\) 的投影，并未要求这些投影与不同复形之间的链映射相容，故不与正文的非自然分裂例子冲突。
\end{solution}
